\documentclass[11pt]{article}
\usepackage[margin=1in]{geometry}
\usepackage{amsmath,amssymb}
\newcommand{\Et}{\mathbb{E}_t}
\newcommand{\Vt}{\operatorname{Var}_t}
\newcommand{\Ct}{\operatorname{Cov}_t}
\newcommand{\wt}{\boldsymbol{\varpi}_t}
\newcommand{\xr}{\mathbf{xR}_{t+1}}
\DeclareMathOperator*{\argmax}{arg\,max}
\title{Problem Set 3\\\large BUSFIN 8200: Asset Pricing Foundations}
\author{}
\date{}
\begin{document}
\setlength{\abovedisplayskip}{6pt}
\setlength{\belowdisplayskip}{6pt}
\begin{center}
{\Large Problem Set 3}\\[3pt]
BUSFIN 8200: Asset Pricing Foundations
\end{center}
\section*{Question 1(a)}
Define $\Sigma_{xR,t}=\Vt[\xr]$. The portfolio choice problem is
\begin{equation*}
\boldsymbol{\varpi}_{p,t}
=\argmax_{\wt}\left\{\Et[R_{p,t+1}]
-\frac{\gamma_t}{2}\Vt[R_{p,t+1}]\right\},
\end{equation*}
subject to
\[
R_{p,t+1}=R_{o,t+1}+\wt'\xr,
\qquad \mathbf{xR}_t=\mathbf{R}_t-R_{o,t}.
\]
Substituting the portfolio return into the objective gives
\begin{align*}
\boldsymbol{\varpi}_{p,t}
&=\argmax_{\wt}\Bigl\{
\Et[R_{o,t+1}+\wt'\xr]
-\frac{\gamma_t}{2}\Vt[R_{o,t+1}+\wt'\xr]\Bigr\}\\
&=\argmax_{\wt}\Bigl\{
\Et[R_{o,t+1}]+\Et[\wt'\xr]\\
&\hspace{3em}-\frac{\gamma_t}{2}\bigl(
\Vt[R_{o,t+1}]+\Vt[\wt'\xr]
+2\Ct(R_{o,t+1},\wt'\xr)\bigr)\Bigr\}.
\end{align*}
Since $\Et[R_{o,t+1}]$ and $\Vt[R_{o,t+1}]$ do not depend on $\wt$,
the maximization can be written as
\begin{align*}
\boldsymbol{\varpi}_{p,t}
&=\argmax_{\wt}\Bigl\{\Et[\wt'\xr]
-\frac{\gamma_t}{2}\Vt[\wt'\xr]
-\gamma_t\Ct(R_{o,t+1},\wt'\xr)\Bigr\}\\
&=\argmax_{\wt}\Bigl\{\wt'\Et[\xr]
-\frac{\gamma_t}{2}\wt'\Vt[\xr]\wt
-\gamma_t\wt'\Ct(R_{o,t+1},\xr)\Bigr\}.
\end{align*}
Taking the derivative with respect to $\wt$ yields
\begin{align*}
\Et[\xr]-\frac{\gamma_t}{2}\,2\Vt[\xr]\wt
-\gamma_t\Ct(R_{o,t+1},\xr)&=0,\\
\Et[\xr]-\gamma_t\Vt[\xr]\wt
-\gamma_t\Ct(R_{o,t+1},\xr)&=0,\\
\gamma_t\Vt[\xr]\wt
&=\Et[\xr]-\gamma_t\Ct(R_{o,t+1},\xr).
\end{align*}
Let $F(\wt)$ denote the objective. Its Hessian is
$H_t=-\gamma_t\Sigma_{xR,t}$. If $\Sigma_{xR,t}$ is positive definite,
then for every $\mathbf d\ne\mathbf 0$,
$\mathbf d'\Sigma_{xR,t}\mathbf d>0$. Since $\gamma_t>0$, it follows that
\[
\mathbf d'H_t\mathbf d=-\gamma_t\mathbf d'\Sigma_{xR,t}\mathbf d<0.
\]
Thus the objective is strictly concave, and the stationary point
$\wt^*$ satisfying $\nabla F(\wt^*)=0$ is the global maximum.
Therefore,
\begin{align*}
\boldsymbol{\varpi}_{p,t}
&=\gamma_t^{-1}\Vt[\xr]^{-1}\Et[\xr]
-\bigl(\gamma_t\Vt[\xr]\bigr)^{-1}
\gamma_t\Ct(R_{o,t+1},\xr)\\
&=\frac{1}{\gamma_t}\Sigma_{xR,t}^{-1}\Et[\xr]
-\Sigma_{xR,t}^{-1}\Ct(R_{o,t+1},\xr).
\end{align*}
If a risk-free asset exists with $R_{o,t+1}=R_{f,t}$, then $\Ct(R_{f,t},\xr)=0$, so
\[
\boxed{\boldsymbol{\varpi}_{p,t}
=\frac{1}{\gamma_t}\Sigma_{xR,t}^{-1}\Et[\xr].}
\]
\clearpage
\section*{Question 1(b)}
From Equation (1.2),
\[
\boldsymbol\varpi_{\gamma,t}
=\frac1{\gamma_t}\Sigma_{xR,t}^{-1}\Et[\mathbf{xR}]
-\Sigma_{xR,t}^{-1}\Ct(R_o,\mathbf{xR}).
\]
Multiplying both sides by $\Sigma_{xR,t}$ gives
\begin{align*}
\Sigma_{xR,t}\boldsymbol\varpi_{\gamma,t}
&=\frac1{\gamma_t}\Sigma_{xR,t}\Sigma_{xR,t}^{-1}\Et[\mathbf{xR}]
-\Sigma_{xR,t}\Sigma_{xR,t}^{-1}\Ct(R_o,\mathbf{xR})\\
&=\frac1{\gamma_t}\Et[\mathbf{xR}]-\Ct(R_o,\mathbf{xR}).
\end{align*}
Therefore,
\[
\Et[\mathbf{xR}]=\gamma_t\bigl(
\Sigma_{xR,t}\boldsymbol\varpi_{\gamma,t}+\Ct(R_o,\mathbf{xR})\bigr).
\]
Since market clearing implies
$\boldsymbol\varpi_{\gamma,t}=\boldsymbol\varpi_{m,t}$,
\[
\Et[\mathbf{xR}]=\gamma_t\bigl(
\Sigma_{xR,t}\boldsymbol\varpi_{m,t}+\Ct(R_o,\mathbf{xR})\bigr).
\]
For the $j$th element,
\[
\Et[xR_j]=\gamma_t\bigl[
\Ct(xR_j,\mathbf{xR})\boldsymbol\varpi_{m,t}+\Ct(R_o,xR_j)\bigr].
\]
Since $R_{w,t+1}=R_{o,t+1}+\boldsymbol\varpi_{m,t}'\xr$,
\begin{align*}
\Et[xR_j]
&=\gamma_t\bigl[\Ct(xR_j,\mathbf{xR})\boldsymbol\varpi_{m,t}
+\Ct(R_o,xR_j)\bigr]\\
&=\gamma_t\Ct(xR_j,R_{w,t+1}).
\end{align*}
Thus,
\[
\boxed{\Et[xR_j]=\gamma_t\Ct(R_w,xR_j),\qquad j=1,\ldots,J.}
\]
\subsection*{Notes on the matrix and covariance expressions}
The conditional mean and market weights are column vectors:
\[
\Et[\mathbf{xR}]=\begin{pmatrix}\Et[xR_1]\\\Et[xR_2]\\\vdots\\\Et[xR_J]\end{pmatrix},
\qquad
\boldsymbol\varpi_{m,t}=\begin{pmatrix}\varpi_{m,1,t}\\\varpi_{m,2,t}\\\vdots\\\varpi_{m,J,t}\end{pmatrix}.
\]
The covariance matrix is
\[
\Sigma_{xR,t}=\begin{pmatrix}
\Vt[xR_1]&\Ct(xR_1,xR_2)&\cdots&\Ct(xR_1,xR_J)\\
\Ct(xR_2,xR_1)&\Vt[xR_2]&\cdots&\Ct(xR_2,xR_J)\\
\vdots&\vdots&\ddots&\vdots\\
\Ct(xR_J,xR_1)&\cdots&\cdots&\Vt[xR_J]
\end{pmatrix}.
\]
\clearpage
For example, its second row gives
\begin{align*}
[\Sigma_{xR,t}\boldsymbol\varpi_{m,t}]_2
&=\Ct(xR_2,xR_1)\varpi_{m,1,t}
+\Vt[xR_2]\varpi_{m,2,t}\\
&\quad+\cdots+\Ct(xR_2,xR_J)\varpi_{m,J,t}\\
&=\Ct(xR_2,\mathbf{xR})\boldsymbol\varpi_{m,t}.
\end{align*}
Here $\mathbf{xR}$ has dimension $J\times1$, while
\[
\Ct(xR_j,\mathbf{xR})=
\begin{pmatrix}\Ct(xR_j,xR_1)&\Ct(xR_j,xR_2)&\cdots&\Ct(xR_j,xR_J)\end{pmatrix}
\]
is a $1\times J$ row vector.
The covariance identity $\operatorname{Cov}(\mathbf X,\mathbf a'\mathbf Y)
=\operatorname{Cov}(\mathbf X,\mathbf Y)\mathbf a$ gives
\[
\Ct(\mathbf{xR},\boldsymbol\varpi_{m,t}'\mathbf{xR})
=\Sigma_{xR,t}\boldsymbol\varpi_{m,t}.
\]
The portfolio return inside the covariance is scalar; the resulting vector
of covariances is a column vector. In particular,
\begin{align*}
\Ct(\mathbf{xR},R_w)
&=\Ct(\mathbf{xR},R_o+\boldsymbol\varpi_{m,t}'\mathbf{xR})\\
&=\Ct(\mathbf{xR},R_o)
+\Ct(\mathbf{xR},\boldsymbol\varpi_{m,t}'\mathbf{xR})\\
&=\Ct(\mathbf{xR},R_o)
+\Ct(\mathbf{xR},\mathbf{xR})\boldsymbol\varpi_{m,t}\\
&=\Ct(\mathbf{xR},R_o)+\Sigma_{xR,t}\boldsymbol\varpi_{m,t}.
\end{align*}
\subsection*{The beta representation}
Since $\Et[xR_j]=\gamma_t\Ct(R_w,xR_j)$ for every $j$, in vector form
\[
\Et[\mathbf{xR}]=\gamma_t\Ct(\mathbf{xR},R_w).
\]
Premultiplying by $\boldsymbol\varpi_{m,t}'$,
\begin{align*}
\boldsymbol\varpi_{m,t}'\Et[\mathbf{xR}]
&=\gamma_t\boldsymbol\varpi_{m,t}'\Ct(\mathbf{xR},R_w),\\
\Et[\boldsymbol\varpi_{m,t}'\mathbf{xR}]
&=\gamma_t\Ct(\boldsymbol\varpi_{m,t}'\mathbf{xR},R_w).
\end{align*}
Since $xR_w=\boldsymbol\varpi_{m,t}'\mathbf{xR}$,
\[
\Et[xR_w]=\gamma_t\Ct(xR_w,R_w).
\]
Combining this with the asset-level equation gives
\[
\Et[xR_j]=\frac{\Ct(R_w,xR_j)}{\Ct(xR_w,R_w)}\Et[xR_w],
\qquad j=1,\ldots,J.
\]
Since $xR_w=R_w-R_o$,
\[
\Ct(xR_w,R_w)=\Ct(R_w-R_o,R_w)
=\Vt[R_w]-\Ct(R_w,R_o).
\]
To obtain Equation (1.5), we additionally require $\Ct(R_w,R_o)=0$,
which holds if $R_o$ is risk-free. Since this condition is not imposed
in the question, Equation (1.5) does not follow in general without an
additional assumption.

Suppose we assume $\Ct(R_w,R_o)=0$. Then
\[
\boxed{\Et[xR_j]=\beta_{j,t}\Et[xR_w],\qquad
\beta_{j,t}=\frac{\Ct(R_w,xR_j)}{\Ct(xR_w,R_w)}.}
\]

\clearpage
\section*{Question 1(c)}
The maximum Sharpe ratio portfolio solves
\[
\boldsymbol{\varpi}_{\max,t}
=\argmax_{\wt}\left\{\operatorname{SR}_t[\wt'\xr]
=\frac{\Et[\wt'\xr]}{\operatorname{Vol}_t[\wt'\xr]}\right\}.
\]
We can write
\[
\operatorname{SR}_t[\wt'\xr]
=\frac{\Et[\wt'\xr]}{\operatorname{Vol}_t[\wt'\xr]}
=\frac{\wt'\Et[\xr]}{\operatorname{Vol}_t[\wt'\xr]}.
\]
Taking the derivative of portfolio volatility gives
\begin{align*}
\frac{\partial\operatorname{Vol}_t[\wt'\xr]}{\partial\wt}
&=\frac{\partial\bigl(\Vt[\wt'\xr]\bigr)^{1/2}}{\partial\wt}
=\frac{\partial(\wt'\Sigma_{xR,t}\wt)^{1/2}}{\partial\wt}\\
&=\frac12(\wt'\Sigma_{xR,t}\wt)^{-1/2}\,2\Sigma_{xR,t}\wt\\
&=\Sigma_{xR,t}\wt(\wt'\Sigma_{xR,t}\wt)^{-1/2}.
\end{align*}
Taking the derivative of the Sharpe ratio with respect to $\wt$ gives
\[
\frac{\Et[\xr]\operatorname{Vol}_t[\wt'\xr]
-\wt'\Et[\xr]\Sigma_{xR,t}\wt
(\wt'\Sigma_{xR,t}\wt)^{-1/2}}
{\wt'\Sigma_{xR,t}\wt}.
\]
Setting $\partial\operatorname{SR}_t/\partial\wt=0$, the notes give
\begin{align*}
\frac{\Et[\xr]\operatorname{Vol}_t[\wt'\xr]}
{\wt'\Sigma_{xR,t}\wt}
&=\frac{\wt'\Et[\xr]\Sigma_{xR,t}\wt}
{(\wt'\Sigma_{xR,t}\wt)\operatorname{Vol}_t[\wt'\xr]},\\
\Et[\xr]\operatorname{Vol}_t[\wt'\xr]
&=\frac{\wt'\Et[\xr]\Sigma_{xR,t}\wt}
{\operatorname{Vol}_t[\wt'\xr]},\\
% Student confirmed these two handwritten expressions use Var, not Vol.
\Et[\xr]\Vt[\wt'\xr]
&=\wt'\Et[\xr]\Sigma_{xR,t}\wt,\\
\Et[\xr]\Vt[\wt'\xr]
&=\Et[xR_{\max}]\Sigma_{xR,t}\wt.
\end{align*}
The notes then give
\[
\wt=\frac{\Vt[\wt'\xr]}{\Et[xR_{\max}]}
\Sigma_{xR,t}^{-1}\Et[\xr],
\qquad
\boldsymbol{\varpi}_{\max,t}
=\frac{\Vt[\boldsymbol{\varpi}_{\max,t}'\xr]}{\Et[xR_{\max}]}
\Sigma_{xR,t}^{-1}\Et[\xr].
\]
Define
\[
c=\frac{\Vt[xR_{\max}]}{\Et[xR_{\max}]},
\qquad \mathbf z=\Sigma_{xR,t}^{-1}\Et[\mathbf{xR}].
\]
Then $\boldsymbol{\varpi}_{\max,t}=c\mathbf z$. Imposing the normalization,
\begin{align*}
\mathbf1'\boldsymbol{\varpi}_{\max,t}=1
&\quad\Longrightarrow\quad
1=\mathbf1'(c\mathbf z)=c\mathbf1'\mathbf z,\\
c&=\frac1{\mathbf1'\mathbf z},\qquad
\boldsymbol{\varpi}_{\max,t}=c\mathbf z
=\frac{\mathbf z}{\mathbf1'\mathbf z}.
\end{align*}
Thus,
\[
\boxed{\boldsymbol{\varpi}_{\max,t}
=\frac{\Sigma_{xR,t}^{-1}\Et[\mathbf{xR}]}
{\mathbf1'\Sigma_{xR,t}^{-1}\Et[\mathbf{xR}]}.}
\]
The normalization implies the full allocation vector
$\mathbf W_{\max,t}=(0,\boldsymbol{\varpi}_{\max,t}')'$.
Consequently, the portfolio returns are
\[
R_{\max,t+1}=\mathbf W_{\max,t}'
\begin{pmatrix}R_{o,t+1}\\\mathbf R_{t+1}\end{pmatrix},
\qquad xR_{\max,t+1}=\boldsymbol{\varpi}_{\max,t}'\xr.
\]
\clearpage
\subsection*{Global maximum}
By the Cauchy--Schwarz inequality,
\[
(\mathbf a'\mathbf b)^2\le(\mathbf a'\mathbf a)(\mathbf b'\mathbf b).
\]
Let
\[
\mathbf a=\Sigma_{xR,t}^{1/2}\wt,
\qquad \mathbf b=\Sigma_{xR,t}^{-1/2}\Et[\mathbf{xR}].
\]
Then
\begin{align*}
(\wt'\Et[\mathbf{xR}])^2
&\le(\wt'\Sigma_{xR,t}\wt)
\bigl(\Et[\mathbf{xR}]'\Sigma_{xR,t}^{-1}\Et[\mathbf{xR}]\bigr),\\
\frac{(\wt'\Et[\mathbf{xR}])^2}{\wt'\Sigma_{xR,t}\wt}
&\le\Et[\mathbf{xR}]'\Sigma_{xR,t}^{-1}\Et[\mathbf{xR}],\\
\operatorname{SR}_t^2(\wt)
&\le\Et[\mathbf{xR}]'\Sigma_{xR,t}^{-1}\Et[\mathbf{xR}],\\
\operatorname{SR}_t(\wt)
&\le\sqrt{\Et[\mathbf{xR}]'\Sigma_{xR,t}^{-1}\Et[\mathbf{xR}]}.
\end{align*}
The upper bound is attained when
\[
\Sigma_{xR,t}^{1/2}\wt
=k\Sigma_{xR,t}^{-1/2}\Et[\mathbf{xR}],
\qquad
\wt=k\Sigma_{xR,t}^{-1}\Et[\mathbf{xR}],\quad k>0.
\]
Hence the portfolio derived from the first-order condition achieves the global
maximum Sharpe ratio, provided that $\Sigma_{xR,t}$ is positive definite and
the normalization admits $k>0$.

\subsection*{Expected excess return and variance}
The expected excess return is
\[
\Et[xR_{\max,t+1}]
=\Et\left[\left(
\frac{\Sigma_{xR,t}^{-1}\Et[\mathbf{xR}]}
{\mathbf1'\Sigma_{xR,t}^{-1}\Et[\mathbf{xR}]}
\right)'\xr\right].
\]
Since $\mathbf1'\Sigma_{xR,t}^{-1}\Et[\mathbf{xR}]$ is a scalar and
$\Sigma_{xR,t}$ is symmetric,
\[
\left(\frac{\Sigma_{xR,t}^{-1}\Et[\mathbf{xR}]}
{\mathbf1'\Sigma_{xR,t}^{-1}\Et[\mathbf{xR}]}\right)'
=\frac{\Et[\mathbf{xR}]'\Sigma_{xR,t}^{-1}}
{\mathbf1'\Sigma_{xR,t}^{-1}\Et[\mathbf{xR}]}.
\]
Therefore,
\begin{align*}
\Et[xR_{\max,t+1}]
&=\Et\left[\frac{\Et[\mathbf{xR}]'\Sigma_{xR,t}^{-1}}
{\mathbf1'\Sigma_{xR,t}^{-1}\Et[\mathbf{xR}]}\xr\right]\\
&=\frac{\Et[\mathbf{xR}]'\Sigma_{xR,t}^{-1}\Et[\mathbf{xR}]}
{\mathbf1'\Sigma_{xR,t}^{-1}\Et[\mathbf{xR}]}.
\end{align*}
The conditional variance is
\begin{align*}
\Vt[xR_{\max}]
&=\Vt[\boldsymbol{\varpi}_{\max,t}'\mathbf{xR}]
=\boldsymbol{\varpi}_{\max,t}'\Vt[\mathbf{xR}]
\boldsymbol{\varpi}_{\max,t}\\
&=\boldsymbol{\varpi}_{\max,t}'\Sigma_{xR,t}
\boldsymbol{\varpi}_{\max,t}\\
&=\frac{\Et[\mathbf{xR}]'\Sigma_{xR,t}^{-1}}
{\mathbf1'\Sigma_{xR,t}^{-1}\Et[\mathbf{xR}]}
\Sigma_{xR,t}
\frac{\Sigma_{xR,t}^{-1}\Et[\mathbf{xR}]}
{\mathbf1'\Sigma_{xR,t}^{-1}\Et[\mathbf{xR}]}\\
&=\frac{\Et[\mathbf{xR}]'\Sigma_{xR,t}^{-1}\Et[\mathbf{xR}]}
{\bigl(\mathbf1'\Sigma_{xR,t}^{-1}\Et[\mathbf{xR}]\bigr)^2}.
\end{align*}
\subsection*{Risk aversion}
Assume $R_{o,t+1}=R_{f,t}$. From Question 1(a),
\[
\wt^*=\frac1{\gamma_t}\Sigma_{xR,t}^{-1}\Et[\mathbf{xR}].
\]
Setting $\boldsymbol{\varpi}_{\max,t}=\wt^*$ gives
\[
\frac{\Sigma_{xR,t}^{-1}\Et[\mathbf{xR}]}
{\mathbf1'\Sigma_{xR,t}^{-1}\Et[\mathbf{xR}]}
=\frac1{\gamma_t}\Sigma_{xR,t}^{-1}\Et[\mathbf{xR}],
\qquad
\gamma_t=\mathbf1'\Sigma_{xR,t}^{-1}\Et[\mathbf{xR}].
\]
Using the mean and variance above,
\begin{align*}
\frac{\Et[xR_{\max}]}{\Vt[xR_{\max}]}
&=\frac{\Et[\mathbf{xR}]'\Sigma_{xR,t}^{-1}\Et[\mathbf{xR}]}
{\mathbf1'\Sigma_{xR,t}^{-1}\Et[\mathbf{xR}]}
\frac{\bigl(\mathbf1'\Sigma_{xR,t}^{-1}\Et[\mathbf{xR}]\bigr)^2}
{\Et[\mathbf{xR}]'\Sigma_{xR,t}^{-1}\Et[\mathbf{xR}]}\\
&=\mathbf1'\Sigma_{xR,t}^{-1}\Et[\mathbf{xR}].
\end{align*}
Thus, assuming $R_{o,t+1}=R_{f,t}$,
\[
\boxed{\gamma_t=\frac{\Et[xR_{\max}]}{\Vt[xR_{\max}]}.}
\]

\clearpage
\section*{Question 1(d)}
\subsection*{Market hedging strategies}
The market index model and hedging strategy are
\begin{align*}
xR_{j,t+1}&=\alpha_{j,t}+\beta_{j,t}xR_{m,t+1}+\epsilon_{j,t+1},\\
xR^H_{j,t+1}&=xR_{j,t+1}-\beta_{j,t}xR_{m,t+1}.
\end{align*}
Substituting the first equation into the second gives
\begin{align*}
xR^H_{j,t+1}
&=\alpha_{j,t}+\beta_{j,t}xR_{m,t+1}
 +\epsilon_{j,t+1}-\beta_{j,t}xR_{m,t+1}\\
&=\alpha_{j,t}+\epsilon_{j,t+1}.
\end{align*}
Thus,
\[
\Et[xR^H_{j,t+1}]=\alpha_{j,t}+\Et[\epsilon_{j,t+1}].
\]
Given $\Et[xR^H_{j,t+1}]=\alpha_{j,t}$, we have
$\Et[\epsilon_{j,t+1}]=0$. Also,
\begin{align*}
\Vt[xR^H_{j,t+1}]
&=\Vt[\alpha_{j,t}+\epsilon_{j,t+1}]
=\Vt[\epsilon_{j,t+1}]=\sigma^2_{\epsilon,j,t},\\
\Ct(xR_m,xR^H_j)
&=\Ct(xR_m,\alpha_{j,t}+\epsilon_{j,t+1})
=\Ct(xR_m,\epsilon_{j,t+1})=0,\\
\Ct(xR^H_j,xR^H_k)
&=\Ct(\alpha_{j,t}+\epsilon_{j,t+1},
\alpha_{k,t}+\epsilon_{k,t+1})\\
&=\Ct(\epsilon_{j,t+1},\epsilon_{k,t+1})=0,
\qquad j\ne k.
\end{align*}
Define the new vector and its conditional mean:
\[
\mathbf X_{t+1}=
\begin{pmatrix}xR_{m,t+1}\\xR^H_{1,t+1}\\\vdots\\xR^H_{J,t+1}\end{pmatrix},
\qquad
\boldsymbol\mu_t=\Et[\mathbf X_{t+1}]
=\begin{pmatrix}\mu_{m,t}\\\alpha_{1,t}\\\vdots\\\alpha_{J,t}\end{pmatrix}.
\]
Since the off-diagonal covariances are zero,
\[
\Sigma_t=\begin{pmatrix}
\sigma^2_{m,t}&0&\cdots&0\\
0&\sigma^2_{\epsilon,1,t}&\cdots&0\\
\vdots&\vdots&\ddots&\vdots\\
0&0&\cdots&\sigma^2_{\epsilon,J,t}
\end{pmatrix},\qquad
\Sigma_t^{-1}=\begin{pmatrix}
1/\sigma^2_{m,t}&0&\cdots&0\\
0&1/\sigma^2_{\epsilon,1,t}&\cdots&0\\
\vdots&\vdots&\ddots&\vdots\\
0&0&\cdots&1/\sigma^2_{\epsilon,J,t}
\end{pmatrix}.
\]
\clearpage
\subsection*{Optimal allocation across the market and hedging strategies}

Assume $\sigma^2_{m,t}>0$ and $\sigma^2_{\epsilon,j,t}>0$ for every $j$,
so that $\Sigma_t$ is positive definite and invertible. Also assume
\[
\mathbf1'\Sigma_t^{-1}\boldsymbol\mu_t
=\frac{\mu_{m,t}}{\sigma^2_{m,t}}
+\sum_{j=1}^J\frac{\alpha_{j,t}}{\sigma^2_{\epsilon,j,t}}>0,
\]
so that normalization preserves the positive scaling of the
Sharpe-ratio-maximizing direction.

From Question 1(c),
\[
\boldsymbol\omega_{\max,t}
=\frac{\Sigma_t^{-1}\boldsymbol\mu_t}
{\mathbf1'\Sigma_t^{-1}\boldsymbol\mu_t}.
\]
Here,
\[
\Sigma_t^{-1}\boldsymbol\mu_t
=\begin{pmatrix}
\mu_{m,t}/\sigma^2_{m,t}\\
\alpha_{1,t}/\sigma^2_{\epsilon,1,t}\\
\vdots\\\alpha_{J,t}/\sigma^2_{\epsilon,J,t}
\end{pmatrix},\qquad
\mathbf1'\Sigma_t^{-1}\boldsymbol\mu_t
=\frac{\mu_{m,t}}{\sigma^2_{m,t}}
+\sum_{k=1}^J\frac{\alpha_{k,t}}{\sigma^2_{\epsilon,k,t}}.
\]
Writing $\boldsymbol\omega_{\max,t}
=(\omega_{m,t},\omega_{1,t},\ldots,\omega_{J,t})'$, we obtain
\begin{align*}
xR_{\max,t+1}
&=\boldsymbol\omega_{\max,t}'\mathbf X_{t+1}\\
&=\omega_{m,t}xR_{m,t+1}
+\sum_{j=1}^J\omega_{j,t}xR^H_{j,t+1}.
\end{align*}
Since $xR_{m,t+1}=R_{m,t+1}-R_{o,t+1}$ and
$xR^H_{j,t+1}=R^H_{j,t+1}-R_{o,t+1}$,
\begin{align*}
xR_{\max,t+1}
&=\omega_{m,t}(R_{m,t+1}-R_{o,t+1})
+\sum_{j=1}^J\omega_{j,t}(R^H_{j,t+1}-R_{o,t+1})\\
&=\omega_{m,t}R_{m,t+1}+\sum_{j=1}^J\omega_{j,t}R^H_{j,t+1}
-\left(\omega_{m,t}+\sum_{j=1}^J\omega_{j,t}\right)R_{o,t+1}.
\end{align*}
The normalization gives
\[
\mathbf1'\boldsymbol\omega_{\max,t}
=\frac{\mathbf1'\Sigma_t^{-1}\boldsymbol\mu_t}
{\mathbf1'\Sigma_t^{-1}\boldsymbol\mu_t}=1,
\qquad \omega_{m,t}+\sum_{j=1}^J\omega_{j,t}=1.
\]
Therefore,
\begin{align*}
xR_{\max,t+1}
&=\omega_{m,t}R_{m,t+1}+\sum_{j=1}^J\omega_{j,t}R^H_{j,t+1}-R_{o,t+1},\\
R_{\max,t+1}
&=\omega_{m,t}R_{m,t+1}+\sum_{j=1}^J\omega_{j,t}R^H_{j,t+1},
\end{align*}
where
\[
\boxed{\omega_{m,t}=
\frac{\mu_{m,t}/\sigma^2_{m,t}}
{\mu_{m,t}/\sigma^2_{m,t}+\sum_{k=1}^J\alpha_{k,t}/\sigma^2_{\epsilon,k,t}},
\qquad
\omega_{j,t}=
\frac{\alpha_{j,t}/\sigma^2_{\epsilon,j,t}}
{\mu_{m,t}/\sigma^2_{m,t}+\sum_{k=1}^J\alpha_{k,t}/\sigma^2_{\epsilon,k,t}}.}
\]
\clearpage
\subsection*{A single hedging portfolio}
For the single hedging portfolio,
\[
xR^H_{t+1}=\sum_{j=1}^J\omega^H_{j,t}xR^H_{j,t+1}.
\]
Its conditional mean is
\begin{align*}
\Et[xR^H]
&=\Et\left[\sum_{j=1}^J\omega^H_{j,t}xR^H_{j,t+1}\right]
=\sum_{j=1}^J\omega^H_{j,t}\Et[xR^H_{j,t+1}],\\
\alpha_{H,t}&=\sum_{j=1}^J\omega^H_{j,t}\alpha_{j,t}.
\end{align*}
Its conditional variance is
\begin{align*}
\Vt[xR^H]
&=\sum_{j=1}^J(\omega^H_{j,t})^2\Vt[xR^H_{j,t+1}]\\
&\quad+2\sum_{j<k}\omega^H_{j,t}\omega^H_{k,t}
\Ct(xR^H_{j,t+1},xR^H_{k,t+1}).
\end{align*}
Since $xR^H_{j,t+1}=\alpha_{j,t}+\epsilon_{j,t+1}$, for $j\ne k$,
\begin{align*}
\Ct(xR^H_{j,t+1},xR^H_{k,t+1})
&=\Ct(\alpha_{j,t}+\epsilon_{j,t+1},\alpha_{k,t}+\epsilon_{k,t+1})\\
&=\Ct(\epsilon_{j,t+1},\epsilon_{k,t+1})=0.
\end{align*}
Thus,
\[
\Vt[xR^H]=\sum_{j=1}^J(\omega^H_{j,t})^2\sigma^2_{\epsilon,j,t}.
\]
\clearpage
\subsection*{Self-note: variance of a weighted sum}
For $z=\sum_{j=1}^J a_jX_j$,
\begin{align*}
\operatorname{Var}(z)
&=\mathbb E[(z-\mathbb E[z])^2]\\
&=\mathbb E\left[\left(\sum_j a_jX_j-\mathbb E\left[\sum_j a_jX_j\right]\right)^2\right]\\
&=\mathbb E\left[\left(\sum_j a_j(X_j-\mathbb E[X_j])\right)^2\right].
\end{align*}
Using $(\sum_j b_j)^2=\sum_j\sum_k b_jb_k$ and
$\operatorname{Cov}(X_j,X_k)=\mathbb E[(X_j-\mathbb E[X_j])(X_k-\mathbb E[X_k])]$,
\begin{align*}
\operatorname{Var}(z)
&=\mathbb E\left[\sum_j\sum_k a_ja_k
(X_j-\mathbb E[X_j])(X_k-\mathbb E[X_k])\right]\\
&=\sum_j\sum_k a_ja_k\operatorname{Cov}(X_j,X_k)\\
&=\sum_j a_j^2\operatorname{Var}(X_j)
+2\sum_{j<k}a_ja_k\operatorname{Cov}(X_j,X_k).
\end{align*}
The last step uses covariance symmetry and separates $j=k$ from $j\ne k$.
\clearpage
\subsection*{Combining the market with the single hedging portfolio}
Write
\[
R_{\max,t+1}=(1-\pi_{H,t})R_{m,t+1}+\pi_{H,t}R^H_{t+1},
\qquad
\mathbf X_{t+1}=\begin{pmatrix}xR_{m,t+1}\\xR^H_{t+1}\end{pmatrix},
\qquad
\boldsymbol\mu_t=\begin{pmatrix}\mu_{m,t}\\\alpha_{H,t}\end{pmatrix}.
\]
Since $xR^H_{j,t+1}=\alpha_{j,t}+\epsilon_{j,t+1}$ and
$\Ct(xR_{m,t+1},\epsilon_{j,t+1})=0$,
\begin{align*}
\Ct(xR_{m,t+1},xR^H_{j,t+1})
&=\Ct(xR_{m,t+1},\alpha_{j,t}+\epsilon_{j,t+1})\\
&=\Ct(xR_{m,t+1},\alpha_{j,t})
+\Ct(xR_{m,t+1},\epsilon_{j,t+1})=0.
\end{align*}
It follows that
\begin{align*}
\Ct(xR_{m,t+1},xR^H_{t+1})
&=\Ct\left(xR_{m,t+1},\sum_{j=1}^J\omega^H_{j,t}xR^H_{j,t+1}\right)\\
&=\sum_{j=1}^J\omega^H_{j,t}\Ct(xR_{m,t+1},xR^H_{j,t+1})=0.
\end{align*}
Therefore,
\[
\Sigma_t=\begin{pmatrix}\sigma^2_{m,t}&0\\0&\sigma^2_{\epsilon,H,t}\end{pmatrix},
\qquad
\Sigma_t^{-1}=\begin{pmatrix}1/\sigma^2_{m,t}&0\\0&1/\sigma^2_{\epsilon,H,t}\end{pmatrix}.
\]
Using Question 1(c),
\begin{align*}
\Sigma_t^{-1}\boldsymbol\mu_t
&=\begin{pmatrix}\mu_{m,t}/\sigma^2_{m,t}\\\alpha_{H,t}/\sigma^2_{\epsilon,H,t}\end{pmatrix},\\
\mathbf1'\Sigma_t^{-1}\boldsymbol\mu_t
&=\mu_{m,t}/\sigma^2_{m,t}+\alpha_{H,t}/\sigma^2_{\epsilon,H,t},\\
\boldsymbol\varpi_{\max,t}
&=\frac{\Sigma_t^{-1}\boldsymbol\mu_t}{\mathbf1'\Sigma_t^{-1}\boldsymbol\mu_t}
=\frac{\begin{pmatrix}\mu_{m,t}/\sigma^2_{m,t}\\\alpha_{H,t}/\sigma^2_{\epsilon,H,t}\end{pmatrix}}
{\mu_{m,t}/\sigma^2_{m,t}+\alpha_{H,t}/\sigma^2_{\epsilon,H,t}}
=\begin{pmatrix}1-\pi_{H,t}\\\pi_{H,t}\end{pmatrix}.
\end{align*}
Thus,
\[
\pi_{H,t}=\frac{\alpha_{H,t}/\sigma^2_{\epsilon,H,t}}
{\mu_{m,t}/\sigma^2_{m,t}+\alpha_{H,t}/\sigma^2_{\epsilon,H,t}}.
\]
Since $xR_{\max,t+1}=(1-\pi_{H,t})xR_{m,t+1}+\pi_{H,t}xR^H_{t+1}$,
\begin{align*}
\Et[xR_{\max,t+1}]&=(1-\pi_{H,t})\mu_{m,t}+\pi_{H,t}\alpha_{H,t},\\
\Vt[xR_{\max}]&=(1-\pi_{H,t})^2\sigma^2_{m,t}
+\pi_{H,t}^2\sigma^2_{\epsilon,H,t}.
\end{align*}
\clearpage
\subsection*{Sharpe-ratio improvement}
The squared Sharpe ratios are
\begin{align*}
\operatorname{SR}_t[xR_{\max}]^2
&=\frac{\Et[xR_{\max}]^2}{\Vt[xR_{\max}]}
=\frac{[(1-\pi_{H,t})\mu_{m,t}+\pi_{H,t}\alpha_{H,t}]^2}
{(1-\pi_{H,t})^2\sigma^2_{m,t}+\pi_{H,t}^2\sigma^2_{\epsilon,H,t}},\\
\operatorname{SR}_t[xR_m]^2&=\frac{\mu_{m,t}^2}{\sigma^2_{m,t}}.
\end{align*}
% Literal transcription of the initial comparison in the handwritten notes;
% the final comparison below uses the opposite subtraction order.
Initially comparing them as in the notes gives
\[
\operatorname{SR}_t[xR_m]^2-\operatorname{SR}_t[xR_{\max}]^2
=\frac{\mu_{m,t}^2}{\sigma^2_{m,t}}
-\frac{[(1-\pi_{H,t})\mu_{m,t}+\pi_{H,t}\alpha_{H,t}]^2}
{(1-\pi_{H,t})^2\sigma^2_{m,t}+\pi_{H,t}^2\sigma^2_{\epsilon,H,t}}.
\]
Define $A=\mu_{m,t}/\sigma^2_{m,t}$ and
$B=\alpha_{H,t}/\sigma^2_{\epsilon,H,t}$. Then
\[
\pi_{H,t}=\frac{B}{A+B},\qquad 1-\pi_{H,t}=\frac{A}{A+B}.
\]
Substituting these weights,
\begin{align*}
\operatorname{SR}_t[xR_{\max}]^2
&=\frac{\left[\dfrac{A}{A+B}\mu_{m,t}
+\dfrac{B}{A+B}\alpha_{H,t}\right]^2}
{\left(\dfrac{A}{A+B}\right)^2\sigma^2_{m,t}
+\left(\dfrac{B}{A+B}\right)^2\sigma^2_{\epsilon,H,t}}\\
&=\frac{[A\mu_{m,t}+B\alpha_{H,t}]^2}
{A^2\sigma^2_{m,t}+B^2\sigma^2_{\epsilon,H,t}}.
\end{align*}
Moreover,
\begin{align*}
A\mu_{m,t}+B\alpha_{H,t}
&=\frac{\mu_{m,t}^2}{\sigma^2_{m,t}}
+\frac{\alpha_{H,t}^2}{\sigma^2_{\epsilon,H,t}},\\
A^2\sigma^2_{m,t}+B^2\sigma^2_{\epsilon,H,t}
&=\left(\frac{\mu_{m,t}}{\sigma^2_{m,t}}\right)^2\sigma^2_{m,t}
+\left(\frac{\alpha_{H,t}}{\sigma^2_{\epsilon,H,t}}\right)^2\sigma^2_{\epsilon,H,t}\\
&=\frac{\mu_{m,t}^2}{\sigma^2_{m,t}}
+\frac{\alpha_{H,t}^2}{\sigma^2_{\epsilon,H,t}}.
\end{align*}
Thus $A\mu_{m,t}+B\alpha_{H,t}
=A^2\sigma^2_{m,t}+B^2\sigma^2_{\epsilon,H,t}$, and
\begin{align*}
\operatorname{SR}_t[xR_{\max}]^2
&=\frac{(A^2\sigma^2_{m,t}+B^2\sigma^2_{\epsilon,H,t})^2}
{A^2\sigma^2_{m,t}+B^2\sigma^2_{\epsilon,H,t}}\\
&=A^2\sigma^2_{m,t}+B^2\sigma^2_{\epsilon,H,t},\\
\operatorname{SR}_t[xR_{\max}]^2-\operatorname{SR}_t[xR_m]^2
&=\frac{\mu_{m,t}^2}{\sigma^2_{m,t}}
+\frac{\alpha_{H,t}^2}{\sigma^2_{\epsilon,H,t}}
-\frac{\mu_{m,t}^2}{\sigma^2_{m,t}}
=\frac{\alpha_{H,t}^2}{\sigma^2_{\epsilon,H,t}}.
\end{align*}
Since the question assumes $\alpha_{H,t}>0$ and
$\sigma_{\epsilon,H,t}>0$, taking the positive square root gives
\[
\boxed{\frac{\alpha_{H,t}}{\sigma_{\epsilon,H,t}}
=\sqrt{\operatorname{SR}_t[xR_{\max}]^2-\operatorname{SR}_t[xR_m]^2},}
\]
where $R_{\max,t+1}=(1-\pi_{H,t})R_{m,t+1}+\pi_{H,t}R^H_{t+1}$.


Equation (1.14) implies that the maximum squared Sharpe ratio attainable by combining the market portfolio with a market-neutral hedging portfolio is

\[
SR[xR_{\max}]^2
=
SR[xR_m]^2
+
\frac{\alpha_H^2}{\sigma_{\epsilon,H}^2}.
\]

Since the market Sharpe ratio is fixed, maximizing the overall Sharpe ratio is equivalent to first choosing the hedging weights \(\omega_j^H\) to maximize \(\alpha_H/\sigma_{\epsilon,H}\), and then choosing the allocation \(\pi_H\) between the market and the optimal hedging portfolio. If \(\alpha_H\neq0\), the hedging portfolio provides an additional investment opportunity that increases the attainable maximum Sharpe ratio. Thus, the market portfolio alone is not mean-variance efficient.

This result has an important implication for SDF theory. For traded excess returns, the pricing condition is \(E[MxR_j]=0\). A linear SDF can be represented as an affine function of the maximum Sharpe ratio portfolio. If the market portfolio were sufficient to span the relevant pricing component of the SDF, a market-only linear specification \(M=a-b_mxR_m\) would imply zero expected excess returns for all market-neutral strategies, because such strategies are uncorrelated with the market and therefore with the SDF. Hence, the pricing condition would require \(\alpha_j=0\).

However, the existence of market-neutral strategies with nonzero alpha implies that the market-only linear SDF is insufficient. Since the maximum Sharpe ratio portfolio may include both the market and additional hedging strategies, the corresponding linear SDF can be expressed as \(M=a-b_mxR_m-\boldsymbol b_F'\boldsymbol F\), where \(\boldsymbol F\) denotes additional traded return factors.

Therefore, reduced-form multifactor models are motivated by the need to capture investment opportunities and priced return variation not explained by the market factor alone. Including additional factors can improve the model's ability to price the cross-section of asset returns, although the portfolio-based motivation does not, by itself, provide a structural economic explanation for the factor risk premia.
\end{document}
